% ---------------------------------------------------------------------
% TÓM TẮT
% ---------------------------------------------------------------------
\thispagestyle{empty}
\begin{center}{\Large\bfseries TÓM TẮT}\end{center}

\noindent\textbf{Tên đề tài:} NGHIÊN CỨU VÀ XÂY DỰNG NỀN TẢNG EDGE AI TÍCH HỢP
AI AGENT VÀ TỐI ƯU MÔ HÌNH YOLO

\noindent\textbf{Sinh viên thực hiện:} Đỗ Quang Hợp

\noindent\textbf{Mã SV:} 68CNCS \hspace{3cm} \textbf{Lớp:} 68CS2

Đồ án nghiên cứu đồng thời hai vấn đề: tối ưu mô hình YOLO cho thiết bị Edge AI
và xây dựng nền tảng quản lý vòng đời AI trên một đội thiết bị phân tán. Ở phần
mô hình, đồ án xây dựng YOLO baseline, khảo sát structured pruning với nhiều tỷ
lệ, Quantization-Aware Training và pipeline kết hợp pruning--fine-tuning--QAT.
Các checkpoint được chuyển đổi sang ONNX và TensorRT ở các chế độ FP16, INT8,
sau đó đánh giá trên Jetson Nano và Jetson Orin Nano theo độ chính xác, latency,
FPS, kích thước model, GPU memory, power và nhiệt độ.

Ở phần nền tảng, đồ án tổ chức hệ thống theo kiến trúc Cloud Control Plane--Edge
Agent--Edge Runtime. Control Plane cung cấp Fleet Management, Model Registry,
Deployment, Command, Monitoring, Incident và Knowledge; Edge Agent thực hiện
đồng bộ desired state, tải và kiểm tra artifact, điều khiển runtime, thu telemetry
và rollback khi triển khai thất bại. AI Copilot kết hợp RAG với tool calling/MCP
để truy vấn tài liệu và trạng thái vận hành. Các hành động thay đổi hệ thống phải
được kiểm tra quyền, lập kế hoạch, phê duyệt và ghi audit.

Giải pháp được đánh giá bằng một kịch bản xuyên suốt từ huấn luyện, tối ưu,
đóng gói và phát hành model đến triển khai trên hai lớp Jetson, theo dõi suy giảm
FPS, tạo incident, phân tích bằng Copilot và rollback. Kết quả định lượng cuối
cùng được để ở trạng thái cần bổ sung cho đến khi hoàn thành phép đo trên thiết
bị thực tế. Đóng góp của đồ án là kết nối nghiên cứu tối ưu YOLO với cơ chế vận
hành model có phiên bản, có khả năng quan sát, khôi phục và hỗ trợ ra quyết định.
\newpage